\section{The complete experiment and its consequences}
\label{sec:critical-law}

\subsection{Uniform kernels and the rectangular reference law}

Let $U$ be a uniform $(D+b)\times D$ matrix. Counting its kernel and then
an injective map from the quotient gives, for
$\max(0,-b)\le k\le D$,
\begin{equation}
 \Pbb(\dim\ker U=k)
 =q^{-k(k+b)}\frac{g_D g_{D+b}}{g_k g_{D-k}g_{k+b}}.
 \label{eq:finite-nullity}
\end{equation}
For other $k$ the probability is zero. For every fixed $k$, this tends to
$\pi_b(k)$. Moreover
\[
 \E q^{\dim\ker U}=1+(q^D-1)q^{-(D+b)},
\]
by counting vectors in the kernel. This bounds the dimension tails and
establishes convergence in total variation, as well as
$\sum_k\pi_b(k)=1$. Conditional on its dimension, the kernel is uniform
over the corresponding Grassmannian, by invariance under $\GL_D(q)$.

We will also use the fixed Gaussian moments
\begin{equation}
 \E_{K\sim\pi_b}\qbinom K r
 =\frac{q^{-r(r+b)}}{g_r}\qquad(r\ge0).
 \label{eq:reference-moments}
\end{equation}
For finite $D$, the left side before passage to the limit is
$\qbinom D r q^{-r(D+b)}$, since each $r$-subspace survives with that
probability. The analogous bound at order $r+1$ ensures uniform
integrability: Gaussian binomials of order $r$ grow at most as $q^{rK}$,
and the order-$r+1$ moment controls that tail. Passing to the limit proves
\eqref{eq:reference-moments}. In particular every fixed exponential
moment $\E q^{rK}$ is finite. The same moments are uniformly bounded for
the actual candidate space, by \cref{ker:aggregate}.

\subsection{Proof of the complete law}

\begin{proof}[Proof of \cref{thm:main}]
Put $b=a-1$. By \cref{ker:theorem}, the actual null candidate space is
within $o(1)$ in total variation of a uniform-row kernel. The latter has
dimension tending to $\pi_b$ and, conditional on dimension $k$, is a
uniform $k$-subspace of $E$. For each fixed $k$, \cref{cor:random-space}
gives $Z\Rightarrow\Bin(m_k,\omega_q)$ in total variation. First truncate
the dimension at a fixed bound, use the finitely many conditional
convergences there, and then let the bound tend to infinity. Tightness of
the dimension proves convergence of $(K,Z)$ to \eqref{eq:Q-limit}.

The exact likelihood is $Z/\lambda_{n,a}$, where
$\lambda_{n,a}=M_np_n^{d+a}\to\lambda_a>0$ by
\eqref{eq:mean-limit}. Since $Z\le m_K$ and
\begin{equation}
 m_K^2=m_K+q(q+1)\qbinom K2,
 \label{eq:second-point-moment}
\end{equation}
the order-one and order-two aggregate inclusion bounds give a uniform
bound on $\E_Q Z^2$. Therefore the size-bias weights are uniformly
integrable. Truncating $Z$ at a fixed bound and using total variation on
the bounded part proves convergence of the planted joint law to
$z\mathsf Q_a(k,z)/\lambda_a$.

The limit has mean
\[
 \E_{\mathsf Q_a}Z
 =\omega_q\E_{\pi_{a-1}}m_K
 =\frac{\omega_qq^{1-a}}{q-1}=\lambda_a.
\]
Thus this size bias is a probability law. Finally,
\begin{equation}
 \frac{m_k\omega_q}{\lambda_a}\pi_{a-1}(k)
 =q^{a-1}(q^k-1)\pi_{a-1}(k)
 =\pi_a(k-1).
 \label{eq:dimension-sizebias}
\end{equation}
Size-biasing a $\Bin(m_k,\omega_q)$ variable yields
$1+\Bin(m_k-1,\omega_q)$, which proves \eqref{eq:P-limit}.
The finite likelihood factorization in \cref{lem:likelihood} establishes
the claim about tests using all observations.
\end{proof}

\subsection{Optimal testing, recovery, and the strict information gain}

\begin{proof}[Proof of \cref{cor:curves}]
The Neyman--Pearson likelihood rule declares the planted law when
$Z>\lambda_{n,a}$; ties may be decided arbitrarily. The exact total
variation is
\[
 \frac12\E_Q\left|Z/\lambda_{n,a}-1\right|.
\]
The joint-law convergence and the second-moment bound just used imply
convergence to \eqref{eq:T}. The exact recovery identity
\eqref{eq:finite-recovery} and convergence of $Q(Z>0)$ give
\eqref{eq:R}. Under the planted limit, conditional on $K=k$, the
probability that $Z=1$ is $(1-\omega_q)^{m_k-1}$, giving
\eqref{eq:U}.

For $a\ge0$, \eqref{eq:threshold-below-one} proves that $Z>0$ is
already the exact finite-sample optimal test. In the limit
$0<\lambda_a<1$. Because $Z$ is integer-valued, the likelihood difference
$Z/\lambda_a-1$ is negative only at $Z=0$. Its expectation is zero,
so half its absolute mean equals $\Pbb_{\mathsf Q_a}(Z=0)$. This is
\eqref{eq:nonnegative}, including its recovery identity.
\end{proof}

\begin{proposition}[Dimension loses information at every fixed offset]
\label{prop:strict}
For every prime power $q$ and every fixed integer $a$,
\[
 T_q(a)>T_q^{\rm dim}(a),\qquad R_q(a)>U_q(a)>F_q(a).
\]
In particular, the nondegeneracy constraint supplies positive-probability
unique recovery at every fixed negative offset, where linear uniqueness
has probability zero.
\end{proposition}

\begin{proof}
By \eqref{eq:dimension-sizebias}, the limiting dimension likelihood is
$q^{a-1}(q^k-1)$. If $a\ge1$, it is below one at $k=0$ and at
least one at every $k\ge1$; hence its total variation is
$\pi_{a-1}(0)=g_\infty/g_{a-1}$. If $a\le0$, the common minimum
dimension is $k_0=1-a$. The dimension likelihood is below one only
there, with deficit $q^{a-1}$. Therefore its total variation is
$q^{a-1}\pi_{a-1}(k_0)=q^{a-1}g_\infty/g_{1-a}$. This proves
\eqref{eq:dimension-curve}.

For strictness, choose any sufficiently large fixed $k$ with
$\pi_{a-1}(k)>0$ and $m_k>\lambda_a$. Conditional on this $k$, the
binomial law gives positive probability both to $Z=0$ and to $Z=m_k$.
Thus the full likelihood minus one has both signs in this conditional
class, and
\[
 \E\bigl[|Z/\lambda_a-1|\mid K=k\bigr]
 >\left|\E[Z/\lambda_a-1\mid K=k]\right|.
\]
Averaging gives the strict testing gain.

Linear uniqueness means $K=1$. Under the planted limiting law its
probability is $\pi_a(0)$, exactly \eqref{eq:linear-curve}. There
is positive planted probability of $K\ge2$ and $Z=1$: choose any
fixed $k\ge2$ in the support and use $0<\omega_q<1$. This proves
$U_q(a)>F_q(a)$, including negative offsets. Finally, on $Z\ge2$
the optimal recovery rule still succeeds with probability $1/Z>0$,
whereas uniqueness fails. That event also has positive limiting
probability, so $R_q(a)>U_q(a)$.
\end{proof}

\subsection{The candidate count is not Poisson}

The earlier critical chi-square calculation identifies only a second
moment. It does not determine the distribution; the new law exhibits a
specific higher-order difference.

\begin{corollary}[Matching variance, different third moment]
\label{cor:not-poisson}
For the null limiting count $Z$,
\begin{align}
 \E Z&=\operatorname{Var}(Z)=\lambda_a,\label{eq:mean-var}\\
 \E\bigl[Z(Z-1)(Z-2)\bigr]
 &=\lambda_a^3+(q-1)\omega_q\lambda_a^2.
 \label{eq:third-factorial}
\end{align}
In particular $Z$ is not Poisson for any fixed $q,a$. The corresponding
limiting chi-square divergence is
\begin{equation}
 \chi^2(\mathsf P_a\Vert\mathsf Q_a)=\lambda_a^{-1}=q^a/c_q.
 \label{eq:chi}
\end{equation}
The finite-experiment chi-square divergence has this same limit.
\end{corollary}

\begin{proof}
Write $M=m_K$. The expected number of ordered linearly independent
projective $r$-tuples in the uniform reference kernel is
\begin{equation}
 \frac{q^{r(1-a)}}{(q-1)^r}.
 \label{eq:independent-points}
\end{equation}
Indeed, the number of such tuples in a $D$-space is
$\prod_{j=0}^{r-1}(q^D-q^j)/(q-1)^r$, and all their points survive
with probability $q^{-rt}$. Pass to the limit; moment control justifies
this equivalently from \eqref{eq:reference-moments}.

Every ordered pair of distinct projective points is independent. Thus
$\E M=q^{1-a}/(q-1)$ and
$\E M(M-1)=q^{2(1-a)}/(q-1)^2$. For an ordered distinct triple,
either it is independent or its third point lies on the projective line
through its first two. In the latter case there are precisely $q-1$
choices for the third point. Therefore
\[
 \E(M)_3=\frac{q^{3(1-a)}}{(q-1)^3}
 +(q-1)\frac{q^{2(1-a)}}{(q-1)^2},
\]
where $(x)_r=x(x-1)\cdots(x-r+1)$. Conditional binomial thinning
multiplies the $r$th factorial moment by $\omega_q^r$. This proves
\eqref{eq:mean-var}--\eqref{eq:third-factorial}. Since
$d\mathsf P_a/d\mathsf Q_a=Z/\lambda_a$, its chi-square divergence
is $\operatorname{Var}(Z)/\lambda_a^2=1/\lambda_a$.
For finite experiments, the uniform bound on a third moment of $Z$,
obtained from aggregate inclusion moments through order three, makes
$Z^2$ uniformly integrable. Thus the second moment and chi-square
divergence converge as well.
\end{proof}

\subsection{An exact optimal algorithm with polynomial expected cost}

\begin{proposition}[Efficient recovery in the critical window]
\label{prop:algorithm}
Fix $q$ and $a$. Given bases for the observed two-planes, the exact
finite-sample Bayes optimal detector and a Bayes optimal hidden-form
recovery rule can be computed with $O_{q,a}(n^6)$ expected field operations
under either $Q_{n,d+a}$ or $P_{n,d+a}$. A rational likelihood comparison
has polynomial bit complexity. The algorithm has no polynomial
worst-case guarantee for candidate enumeration.
\end{proposition}

\begin{proof}
For a basis $(u_i,v_i)$ of $\Pi_i$, its equation is
$u_i^{\mathsf T}Gv_i=0$. Form the $t\times D$ coefficient matrix and
compute a basis $G_1,\ldots,G_K$ of $N$ by Gaussian elimination. This
takes $O(n^6)$ field operations because $t,D=O(n^2)$.

Enumerate the $m_K$ projective coefficient vectors, normalized so that
their first nonzero coordinate is one. Form the corresponding matrices
$\sum_jx_jG_j$ and compute their ranks. This costs at most
$O(m_K(Kn^2+n^3))$ field operations and determines $Z$ and all compatible
nondegenerate forms. Compare $Z$ to $M_np_n^t$ for the optimal test.
Return any compatible nondegenerate form for Bayes optimal recovery;
the posterior on that list is uniform.

The fixed-order aggregate moments imply
$\sup_n\E_Q q^{rK}<\infty$ for every fixed $r$. Under the planted law,
\[
 \E_P f(K)=\frac{\E_Q[Zf(K)]}{M_np_n^t},\qquad
 Z\le m_K\le q^K/(q-1),
\]
and the denominator tends to the positive constant $\lambda_a$.
Consequently the order-three null moment bounds
$\E_Q[m_K(1+K)]$ and $\E_P[m_K(1+K)]$ uniformly. The expected
enumeration cost is $O_{q,a}(n^3)$, dominated by elimination.
The likelihood threshold is a rational number built from powers of
integers of $O_q(n)$ bits raised to $t=O(n^2)$, together with the exact
form count. Its numerator and denominator have $O_{q,a}(n^3)$ bits,
so exact comparison is polynomial-time.
\end{proof}

For $a\ge0$ the detector can stop as soon as it finds a nondegenerate
candidate, since existence is the exact optimal test. To certify that the
candidate is unique, the rest of the projective space must be checked.
The distinction between an optimal guess and a uniqueness certificate is
reflected in the separate curves $R_q$ and $U_q$.
